\section{引言}

随着大语言模型（LLM）和视觉语言模型（VLM）的快速发展，研究者越来越关注构建 GUI（Graphical User Interface）智能体，使其能够理解自然语言指令，并在桌面端~\cite{zhang2024ufo, zhang2025ufo2}、移动设备~\cite{wang2024mobileagentautonomousmultimodalmobile} 和 Web 应用~\cite{zheng2024gpt} 等不同平台的软件界面中自主交互。有效的 GUI 智能体需要两项核心能力：（i）多模态感知，用于理解视觉线索和语言线索；（ii）动作执行，用于通过鼠标、键盘或触摸屏与数字环境交互~\cite{zhang2024large, wang2024large}。早期系统依赖结构化元数据（\eg HTML、DOM 树或视图层级）~\cite{zhang2024ufo}，但这类数据在不同平台上往往噪声较多、不一致，甚至不可用。因此，近期工作更强调视觉 GUI 智能体，即直接从渲染后的截图中感知界面，类似人类用户的交互方式~\cite{xu2024aguvis}。在这一范式中，一个核心挑战是\emph{视觉定位}：将自然语言计划映射到屏幕区域。大多数已有方法把它视为坐标生成任务，通过 LLM 使用的同一套文本生成机制输出屏幕位置（\eg ``x=0.125, y=0.23''）~\cite{gou2024navigating}。

然而，用坐标生成来表示 GUI 动作，也就是让模型以文本 token 形式输出屏幕位置（\eg, \texttt{x=...}, \texttt{y=...}），会带来几个内在限制。首先，\textit{空间-语义对齐较弱}：生成离散坐标 token 要求模型通过语言建模头，把视觉输入隐式映射成数字输出，却没有明确的空间归纳偏置。这一过程效率低、数据需求大，并且由于缺少直接把视觉特征连接到动作位置的监督，容易出错。其次，\textit{监督信号具有歧义}：许多 GUI 动作，例如点击按钮内部，允许一组有效目标位置。然而，基于坐标的方法通常把任务当作单点预测，惩罚所有偏离，包括合理偏离，因此无法刻画人类交互中的自然歧义。最后，\textit{视觉特征与动作空间之间存在粒度不匹配}：坐标是连续且高分辨率的，而 Vision Transformer（ViT）~\cite{dosovitskiy2021an} 等视觉模型基于 patch 级特征工作。这种不匹配迫使模型从粗粒度视觉 token 中推断密集的像素级动作，从而削弱其对多样屏幕布局和分辨率的泛化能力。

尽管一些近期方法~\cite{qin2025ui} 尝试通过预测边界框而非单点来丰富空间定位，它们仍然把这些区域表示为原始坐标字符串（\eg \texttt{x\_min}, \texttt{y\_min}, \texttt{x\_max}, \texttt{y\_max}），这些字符串与视觉特征本身是脱节的。如果没有 ROI pooling~\cite{sun2019roi} 或空间注意力机制~\cite{zhu2019empirical} 等架构组件，这类方法仍难以弥合语言意图与视觉动作定位之间的差距。

重新审视人类如何与数字界面交互，可以得到一个关键观察：\textit{人类在行动前并不会计算精确屏幕坐标，而是感知目标元素并直接与其交互}。受此启发，我们提出 \sys：一种在 VLM 上增加注意力式动作头的模型，用于实现更接近人类行为的无坐标视觉定位。不同于把动作定位视为坐标预测任务的已有方法，\sys 学习直接关注相关视觉区域，而不依赖数值坐标生成。\sys 的核心是一个专用的 \texttt{<ACTOR>} token，它通过联合处理视觉输入和自然语言指令来编码定位上下文。随后，注意力机制学习将该 token 与截图中的视觉 patch token 对齐，从而找到最相关的 GUI 区域。得到的注意力图可以自然地识别界面中的可操作区域。

为了处理 GUI 交互中的固有歧义，即一个 UI 元素（\eg 一个按钮）内部可能有多个点都是有效的，\sys 使用多 patch 监督进行训练。所有与真实边界框重叠的视觉 patch 都标为正例，其余 patch 标为负例。这种监督策略让模型能够容忍空间歧义，并减少对合理动作变化的过度惩罚。此外，由于 \sys 直接在视觉骨干的原生空间分辨率上定位动作，它避免了已有方法中的粒度不匹配问题，并能更稳健地泛化到不同屏幕尺寸、分辨率和布局。最后，为了支持决策细化，我们进一步为 GUI-Actor 引入一个轻量级定位验证器，用于评估多个候选区域，并选择最可能用于动作执行的区域。

我们的贡献可以概括如下：
\begin{enumerate}[leftmargin=*, label=\arabic*.]
    \item 我们重新审视了近期用于 GUI 智能体视觉定位的坐标生成方法，指出其局限，包括空间-语义对齐较弱、监督目标有歧义以及特征粒度不匹配，并提出 \sys，一种能够有效缓解这些问题的新型无坐标方法。
    \item 我们设计了一个注意力式动作头，它可以在单次前向传播中生成多个候选区域，从而为搜索策略等下游模块提供灵活性。
    \item 我们引入一个定位验证器，用于从动作注意力图提出的候选区域中选择最可能的动作区域。我们还展示了该验证器可以很容易地接入其他定位方法，并进一步提升性能。
    \item 大量实验表明，在多个 GUI 动作定位基准上，\sys 优于在相近数据规模上训练的最新方法，并且对未见过的屏幕尺寸和分辨率表现出更强鲁棒性。值得注意的是，\sys 的 2B 版本甚至超过了若干 7B 级竞争模型。此外，借助验证器，采用轻量训练的 \sys（\ie 冻结骨干 LLM，仅微调动作头中新引入的约 $\sim$100M 参数）可以有效赋予底层 VLM 定位能力，同时不损害其通用能力。
\end{enumerate}
